% Options for packages loaded elsewhere
\PassOptionsToPackage{unicode}{hyperref}
\PassOptionsToPackage{hyphens}{url}
\PassOptionsToPackage{dvipsnames,svgnames,x11names}{xcolor}
%
\documentclass[
  10pt,
  a4paper,
]{article}
\usepackage{amsmath,amssymb}
\usepackage{iftex}
\ifPDFTeX
  \usepackage[T1]{fontenc}
  \usepackage[utf8]{inputenc}
  \usepackage{textcomp} % provide euro and other symbols
\else % if luatex or xetex
  \usepackage{unicode-math} % this also loads fontspec
  \defaultfontfeatures{Scale=MatchLowercase}
  \defaultfontfeatures[\rmfamily]{Ligatures=TeX,Scale=1}
\fi
\usepackage{lmodern}
\ifPDFTeX\else
  % xetex/luatex font selection
\fi
% Use upquote if available, for straight quotes in verbatim environments
\IfFileExists{upquote.sty}{\usepackage{upquote}}{}
\IfFileExists{microtype.sty}{% use microtype if available
  \usepackage[]{microtype}
  \UseMicrotypeSet[protrusion]{basicmath} % disable protrusion for tt fonts
}{}
\makeatletter
\@ifundefined{KOMAClassName}{% if non-KOMA class
  \IfFileExists{parskip.sty}{%
    \usepackage{parskip}
  }{% else
    \setlength{\parindent}{0pt}
    \setlength{\parskip}{6pt plus 2pt minus 1pt}}
}{% if KOMA class
  \KOMAoptions{parskip=half}}
\makeatother
\usepackage{xcolor}
\usepackage[margin=1in]{geometry}
\usepackage{color}
\usepackage{fancyvrb}
\newcommand{\VerbBar}{|}
\newcommand{\VERB}{\Verb[commandchars=\\\{\}]}
\DefineVerbatimEnvironment{Highlighting}{Verbatim}{commandchars=\\\{\}}
% Add ',fontsize=\small' for more characters per line
\newenvironment{Shaded}{}{}
\newcommand{\AlertTok}[1]{\textcolor[rgb]{1.00,0.00,0.00}{\textbf{#1}}}
\newcommand{\AnnotationTok}[1]{\textcolor[rgb]{0.38,0.63,0.69}{\textbf{\textit{#1}}}}
\newcommand{\AttributeTok}[1]{\textcolor[rgb]{0.49,0.56,0.16}{#1}}
\newcommand{\BaseNTok}[1]{\textcolor[rgb]{0.25,0.63,0.44}{#1}}
\newcommand{\BuiltInTok}[1]{\textcolor[rgb]{0.00,0.50,0.00}{#1}}
\newcommand{\CharTok}[1]{\textcolor[rgb]{0.25,0.44,0.63}{#1}}
\newcommand{\CommentTok}[1]{\textcolor[rgb]{0.38,0.63,0.69}{\textit{#1}}}
\newcommand{\CommentVarTok}[1]{\textcolor[rgb]{0.38,0.63,0.69}{\textbf{\textit{#1}}}}
\newcommand{\ConstantTok}[1]{\textcolor[rgb]{0.53,0.00,0.00}{#1}}
\newcommand{\ControlFlowTok}[1]{\textcolor[rgb]{0.00,0.44,0.13}{\textbf{#1}}}
\newcommand{\DataTypeTok}[1]{\textcolor[rgb]{0.56,0.13,0.00}{#1}}
\newcommand{\DecValTok}[1]{\textcolor[rgb]{0.25,0.63,0.44}{#1}}
\newcommand{\DocumentationTok}[1]{\textcolor[rgb]{0.73,0.13,0.13}{\textit{#1}}}
\newcommand{\ErrorTok}[1]{\textcolor[rgb]{1.00,0.00,0.00}{\textbf{#1}}}
\newcommand{\ExtensionTok}[1]{#1}
\newcommand{\FloatTok}[1]{\textcolor[rgb]{0.25,0.63,0.44}{#1}}
\newcommand{\FunctionTok}[1]{\textcolor[rgb]{0.02,0.16,0.49}{#1}}
\newcommand{\ImportTok}[1]{\textcolor[rgb]{0.00,0.50,0.00}{\textbf{#1}}}
\newcommand{\InformationTok}[1]{\textcolor[rgb]{0.38,0.63,0.69}{\textbf{\textit{#1}}}}
\newcommand{\KeywordTok}[1]{\textcolor[rgb]{0.00,0.44,0.13}{\textbf{#1}}}
\newcommand{\NormalTok}[1]{#1}
\newcommand{\OperatorTok}[1]{\textcolor[rgb]{0.40,0.40,0.40}{#1}}
\newcommand{\OtherTok}[1]{\textcolor[rgb]{0.00,0.44,0.13}{#1}}
\newcommand{\PreprocessorTok}[1]{\textcolor[rgb]{0.74,0.48,0.00}{#1}}
\newcommand{\RegionMarkerTok}[1]{#1}
\newcommand{\SpecialCharTok}[1]{\textcolor[rgb]{0.25,0.44,0.63}{#1}}
\newcommand{\SpecialStringTok}[1]{\textcolor[rgb]{0.73,0.40,0.53}{#1}}
\newcommand{\StringTok}[1]{\textcolor[rgb]{0.25,0.44,0.63}{#1}}
\newcommand{\VariableTok}[1]{\textcolor[rgb]{0.10,0.09,0.49}{#1}}
\newcommand{\VerbatimStringTok}[1]{\textcolor[rgb]{0.25,0.44,0.63}{#1}}
\newcommand{\WarningTok}[1]{\textcolor[rgb]{0.38,0.63,0.69}{\textbf{\textit{#1}}}}
\setlength{\emergencystretch}{3em} % prevent overfull lines
\providecommand{\tightlist}{%
  \setlength{\itemsep}{0pt}\setlength{\parskip}{0pt}}
\setcounter{secnumdepth}{-\maxdimen} % remove section numbering
\ifLuaTeX
  \usepackage{selnolig}  % disable illegal ligatures
\fi
\usepackage{bookmark}
\IfFileExists{xurl.sty}{\usepackage{xurl}}{} % add URL line breaks if available
\urlstyle{same}
\hypersetup{
  colorlinks=true,
  linkcolor={blue},
  filecolor={Maroon},
  citecolor={Blue},
  urlcolor={blue},
  pdfcreator={LaTeX via pandoc}}

\author{}
\date{}

\begin{document}

\section{DDTA research work plan after documentation-layer
closure}\label{ddta-research-work-plan-after-documentation-layer-closure}

\textbf{WORK PLAN - REVISION 4}\\
\textbf{Status:} active post-BA0-R execution plan.\\
\textbf{Supersedes:} Revision 3 only for forward execution state; R1-R3
remain historical research records.

\subsection{1. Fixed prior state}\label{1-fixed-prior-state}

\begin{itemize}
\tightlist
\item
  Chapters 2-4: CLOSED / FINAL for current thesis scope.
\item
  Documentation layer: CLOSED.
\item
  BA0-R systems-modeling prior-art research: CLOSED.
\item
  No Chapter 2-4 reopen trigger is introduced by BA0-T1/T2.
\item
  ThreatForge is a case study/tooling instrument, not DDTA semantic
  authority.
\end{itemize}

\subsection{2. Scope correction carried into
R4}\label{2-scope-correction-carried-into-r4}

DDTA evaluation begins from governed documentation intentionally
authored to satisfy an analysis-readiness / portability contract.

Reliable automatic transformation of arbitrary narrative, legacy or
unstructured documentation is \textbf{not} a BA0 requirement and is
outside the primary thesis evaluation. Future NLP/LLM work may assist
migration or candidate extraction but must retain
candidate/review/provenance boundaries.

\subsection{3. Base Analysis responsibility now under
test}\label{3-base-analysis-responsibility-now-under-test}

Base Analysis is the accepted methodology-neutral analytical boundary
that should support, from the same governed project meaning:

\begin{itemize}
\tightlist
\item
  progressive human comprehension and drill-down;
\item
  shared identity/provenance across analyses;
\item
  method-specific projections without core contamination;
\item
  localized diagnostics for missing/ambiguous/inconsistent information;
\item
  governed corrective feedback;
\item
  change-aware revalidation and targeted re-analysis.
\end{itemize}

Semantic canonicality does not imply one permanently materialized graph.

\subsection{4. BA0 trial sequence}\label{4-ba0-trial-sequence}

\subsubsection{BA0-T1 - FR-3.4 application
trial}\label{ba0-t1---fr-34-application-trial}

\textbf{Status: COMPLETED / PROVISIONAL.}

Supported a shared semantic contract and origin distinctions; falsified
\texttt{every\ useful\ semantic\ responsibility\ must\ be\ a\ first-class\ type};
did not prove one persisted canonical graph necessary.

Its historical successor pointer to a conflict/addition stress test is
superseded by the dedicated successor-supersession note in
\texttt{methodology/}.

\subsubsection{BA0-T2 - analysis consumption, teachability and governed
corrective
feedback}\label{ba0-t2---analysis-consumption-teachability-and-governed-corrective-feedback}

\textbf{Status: COMPLETED / PROVISIONAL PASS.}

Required checks:

\begin{enumerate}
\def\labelenumi{\arabic{enumi}.}
\tightlist
\item
  a human can navigate a progressive overview-to-source projection over
  the facial-access branch;
\item
  a concrete bounded analysis consumer can reuse the same accepted
  project semantics without redefining them;
\item
  insufficient/missing information can be returned as a source-targeted
  corrective documentation candidate;
\item
  the trial does not rely on automatic arbitrary-narrative extraction or
  one persisted graph.
\end{enumerate}

The bounded consumer is a pre-overlay STRIDE transfer probe used only to
pressure-test the shared boundary. It is not the formal O2 STRIDE
implementation.

\subsubsection{BA0-T3 - responsibility and non-goals closure
review}\label{ba0-t3---responsibility-and-non-goals-closure-review}

\textbf{NEXT / NOT EXECUTED BY THIS PLAN REVISION.}

Review BA0-R + T1 + T2 together and decide only:

\begin{itemize}
\tightlist
\item
  precise Base Analysis responsibility statement;
\item
  explicit non-goals;
\item
  whether progressive teachability is a required Base Analysis
  capability;
\item
  whether accepted semantic identity may be reused or reproducibly
  rebuilt;
\item
  whether diagnostics and governed corrective feedback belong inside the
  BA0 responsibility boundary;
\item
  whether remaining uncertainty belongs to BA1-BA6 rather than BA0.
\end{itemize}

\textbf{Exit:} BA0 becomes CLOSED only if these responsibilities are
precise enough to begin ontology work without smuggling in a
pre-approved type list.

\subsection{5. Sequence after BA0
closure}\label{5-sequence-after-ba0-closure}

Do not skip phases:

\begin{Shaded}
\begin{Highlighting}[]
\NormalTok{BA1 {-} Minimal BAE ontology}
\NormalTok{BA2 {-} Relations and canonical action vocabulary}
\NormalTok{BA3 {-} Document{-}to{-}BAE derivation, provenance and authority}
\NormalTok{BA4 {-} Multi{-}level human and method projections}
\NormalTok{BA5 {-} Controlled vocabulary and optional authoring/extraction assistance}
\NormalTok{BA6 {-} Base Analysis closure and regression}
\end{Highlighting}
\end{Shaded}

\subsubsection{BA1 - Minimal BAE
ontology}\label{ba1---minimal-bae-ontology}

Start from responsibilities forced jointly by governed DDTA corpora,
BA0-R evidence and BA0 trials. No type from historical
\texttt{primary-analysis-focus} studies, STRIDE, SysML, ThreatForge or
another notation is accepted by default.

\subsubsection{BA2 - Relations and canonical action
vocabulary}\label{ba2---relations-and-canonical-action-vocabulary}

Define the smallest stable semantic relation vocabulary. Keep semantic
relation, canonical lexical label and authoring synonym assistance
separate.

\subsubsection{BA3 - Document-to-BAE derivation, provenance and
authority}\label{ba3---document-to-bae-derivation-provenance-and-authority}

Close identity, source locators, grounded/derived/reviewed-addition
boundaries, acceptance/rejection, stale state and document-to-analysis
traceability. Automatic candidate extraction may be evaluated as
optional assistance but is not required for thesis validity.

\subsubsection{BA4 - Multi-level human and method
projections}\label{ba4---multi-level-human-and-method-projections}

Views are derived from the same accepted Base Analysis semantics,
whether the representation is reused or reproducibly rebuilt. Test both
progressive human teachability and analysis-specific views. No view is
source authority.

\subsubsection{BA5 - Controlled vocabulary and optional
assistance}\label{ba5---controlled-vocabulary-and-optional-assistance}

Vocabulary reuse, duplicate suggestions, LLM/NLP proposals and candidate
extraction are optional support. None may silently establish identity,
equivalence, ownership, project fact or BAE acceptance.

\subsubsection{BA6 - Base Analysis closure and
regression}\label{ba6---base-analysis-closure-and-regression}

Regress closed corpora and at least one structurally different holdout.
Verify semantic identity stability, source traceability, projection
consistency, diagnostic behavior, change/re-analysis scope, method
neutrality and human drill-down support.

\subsection{6. Analysis envelope after Base Analysis
closure}\label{6-analysis-envelope-after-base-analysis-closure}

Only after the Base Analysis gate:

\begin{itemize}
\tightlist
\item
  A1 - AnalysisRecord / execution envelope.
\item
  A2 - common reviewed Finding boundary.
\item
  A3 - change/provenance integration.
\end{itemize}

Then define the generic overlay contract and full methodology
evaluations.

\subsection{7. Method-specific
evaluations}\label{7-method-specific-evaluations}

Planned concrete evaluations remain STRIDE and STRIDE-AI over the same
common core. They are demonstrators of the overlay boundary, not proof
of universal method support.

The BA0-T2 STRIDE probe does not count as the formal STRIDE evaluation;
it exists only to falsify the Base Analysis responsibility before
ontology design.

\subsection{8. Corrective feedback
boundary}\label{8-corrective-feedback-boundary}

The general DDTA loop is broader than
\texttt{Finding\ -\textgreater{}\ SecurityRequirement}:

\begin{Shaded}
\begin{Highlighting}[]
\NormalTok{analysis / validation / diagnostic}
\NormalTok{ {-}\textgreater{} source{-}localized corrective candidate}
\NormalTok{ {-}\textgreater{} governed review}
\NormalTok{ {-}\textgreater{} correction / clarification / supersession / new requirement or decision}
\NormalTok{ {-}\textgreater{} Base Analysis revalidation}
\NormalTok{ {-}\textgreater{} impacted re{-}analysis}
\end{Highlighting}
\end{Shaded}

The thesis may still evaluate the narrower accepted
security-Finding-to-SecurityRequirement path as a specific research
contribution. BA0 must not assume every correction has the same document
type.

\subsection{9. Next authorized
microstep}\label{9-next-authorized-microstep}

Execute only:

\begin{quote}
\textbf{BA0-T3 - responsibility and non-goals closure review.}
\end{quote}

Do not start BA1, formal STRIDE overlay design, Common Finding schema or
implementation work unless BA0 closes explicitly.

\end{document}
